% v2.3.1 figure UID: FIG-P109-01
% Proposed post-v2.3.1 polish for FIG-P109-01: protect key-label size and stroke hierarchy.
\tikzset{
  slfig-FIG-P109-01/.style={font=\fontsize{9.2pt}{11.2pt}\selectfont,
    every path/.append style={line cap=round,line join=round}},
  slfig-FIG-P109-01-label/.style={font=\fontsize{9.2pt}{11.2pt}\selectfont,
    fill=white,fill opacity=.92,text opacity=1,inner sep=1.2pt},
  slfig-FIG-P109-01-note/.style={slfig note,rounded corners=1.4mm,
    align=center,font=\fontsize{9.2pt}{11.2pt}\selectfont,
    inner xsep=6pt,inner ysep=4pt,line width=.55pt}
}
\pgfkeysifdefined{/tikz/alt}{}{\tikzset{alt/.code={}}}
\begin{figure}[H]
\centering
\begin{tikzpicture}[slfig-FIG-P109-01,x=1.35cm,y=1.35cm,
  alt={对象—关系—结论：凸集中任意两点的线段仍位于可行域内}]
  \path[fill=SLBlueBG,draw=SLBlue,line width=.82pt]
    plot[smooth cycle,tension=.82] coordinates {
      (-3.7,-.6) (-2.7,-1.45) (-.5,-1.65) (2.0,-1.25)
      (3.55,-.15) (3.2,1.45) (1.2,2.15) (-1.3,2.05) (-3.25,1.15)};
  \coordinate (x) at (-2.35,-.62);
  \coordinate (y) at (2.42,1.02);
  \draw[draw=SLBlue,line width=1.05pt] (x)--(y);
  \fill[SLInk] (x) circle (2.35pt) node[below left=3mm] {$x$};
  \fill[SLInk] (y) circle (2.35pt) node[above right=3mm] {$y$};
  \foreach \f in {.25,.50,.75}{\fill[SLTeal] ($(x)!\f!(y)$) circle (2.25pt);}
  \node[slfig-FIG-P109-01-label,anchor=south east]
    at ($(x)!.50!(y)+(0,3mm)$) {$z=\lambda x+(1-\lambda)y$};
  \node[font=\fontsize{9.2pt}{11.2pt}\selectfont,text=SLInk!78,fill=white,fill opacity=1,text opacity=1,inner sep=1.2pt,anchor=north east]
    at (3.20,1.86) {凸可行域 $C$};
  \node[slfig-FIG-P109-01-note,anchor=north]
    at (0,-1.92) {$x,y\in C,\ \lambda\in[0,1]\ \Longrightarrow\
      \lambda x+(1-\lambda)y\in C$};
\end{tikzpicture}
\caption{凸集中任意两点的线段仍位于可行域内}\label{fig:V1-C07-convex-set}
\end{figure}
